Each new line crossing the interior increases the number of regions by at least $1$. The three lines therefore produce at least $4$ regions. A new line can cross at most $k-1$ earlier lines, producing at most $k$ new regions at step $k$; thus there are at most $1+1+2+3=7$.

All four values occur. Three parallel lines give $4$. Two parallel lines and a third crossing only one inside the square give $5$. Three concurrent lines give $6$. Three lines with three distinct pairwise intersections inside the square give $7$. Hence the requested sum is $\boxed{4+5+6+7=22}$.
\begin{asy}
texpreamble("\usepackage{amsmath,amssymb}");
size(240);
defaultpen(fontsize(9pt));
pair A=(0,0),B=(6,0),C=(6,6),D=(0,6);draw(A--B--C--D--cycle);label("$A$",A,(-1,-1));label("$B$",B,(1,-1));label("$C$",C,(1,1));label("$D$",D,(-1,1));draw((0,1)--(6,1));draw((0,3)--(6,3));draw((0,5)--(6,5));label("$n=4$",(3,-1));
\end{asy}
\begin{asy}
texpreamble("\usepackage{amsmath,amssymb}");
size(240);
defaultpen(fontsize(9pt));
pair A=(0,0),B=(6,0),C=(6,6),D=(0,6);draw(A--B--C--D--cycle);label("$A$",A,(-1,-1));label("$B$",B,(1,-1));label("$C$",C,(1,1));label("$D$",D,(-1,1));draw((0,2)--(6,2));draw((0,4)--(6,4));draw((0,1)--(6,3));label("$n=5$",(3,-1));
\end{asy}
\begin{asy}
texpreamble("\usepackage{amsmath,amssymb}");
size(240);
defaultpen(fontsize(9pt));
pair A=(0,0),B=(6,0),C=(6,6),D=(0,6);draw(A--B--C--D--cycle);label("$A$",A,(-1,-1));label("$B$",B,(1,-1));label("$C$",C,(1,1));label("$D$",D,(-1,1));draw((0,3)--(6,3));draw((3,0)--(3,6));draw((0,0)--(6,6));label("$n=6$",(3,-1));
\end{asy}
\begin{asy}
texpreamble("\usepackage{amsmath,amssymb}");
size(240);
defaultpen(fontsize(9pt));
pair A=(0,0),B=(6,0),C=(6,6),D=(0,6);draw(A--B--C--D--cycle);label("$A$",A,(-1,-1));label("$B$",B,(1,-1));label("$C$",C,(1,1));label("$D$",D,(-1,1));draw((0,2)--(6,2));draw((1,0)--(4,6));draw((5,0)--(2,6));label("$n=7$",(3,-1));
\end{asy}
